\section{Idea central de Model-Based RL}

\subsection{Qué es Model-Based RL}

A diferencia de los métodos model-free, model-based RL aprende de forma explícita un \textbf{modelo de dinámica} (dynamics model) del entorno y después lo usa para tomar decisiones.

\begin{importantbox}{Modelo de dinámica}
Se aprende una función $f_\phi(\mathbf{s}, \mathbf{a})$ que aproxima la transición de estados del entorno $p(\mathbf{s}' | \mathbf{s}, \mathbf{a})$:
$$f_\phi(\mathbf{s}, \mathbf{a}) \approx \mathbf{s}'$$
El objetivo del entrenamiento es minimizar el error de predicción:
$$\min_\phi \sum_i \|f_\phi(\mathbf{s}_i, \mathbf{a}_i) - \mathbf{s}_i'\|^2$$
donde los datos $\mathcal{D} = \{(\mathbf{s}, \mathbf{a}, \mathbf{s}')_i\}$ provienen de la interacción con el entorno.
\end{importantbox}

\begin{figure}[H]
\centering
\includegraphics[width=0.9\textwidth]{slides-images/slide-05.jpg}
\caption{Temario de esta clase: contenidos centrales de Model-Based RL\protect\footnotemark}
\end{figure}
\footnotetext{Slides, página 5.}

\subsection{Por qué aprender un modelo}

\begin{knowledgebox}{Disyuntiva central entre Model-Based y Model-Free}
\begin{itemize}
    \item \textbf{Eficiencia de datos}: los métodos model-based suelen ser mucho más eficientes en el uso de datos que los métodos model-free. El modelo aprendido puede usarse para «imaginar» un sinnúmero de trayectorias posibles sin necesidad de ejecutarlas en la realidad.
    \item \textbf{Costo}: el modelo no es perfecto, y las decisiones basadas en un modelo inexacto pueden degradar el desempeño. Este es el problema del \textbf{sesgo del modelo} (model bias).
\end{itemize}
\end{knowledgebox}

\subsection{Resumen de la sección}

La idea central de Model-Based RL es «primero aprender un modelo del mundo y después usar ese modelo para planificar». Resulta especialmente valioso cuando los datos son escasos (por ejemplo, en robótica), pero exige manejar con cuidado el error del modelo.

